\documentclass[10pt,a4paper]{article}
\usepackage[utf8]{inputenc}
\usepackage[a4paper,left=3cm,right=2cm]{geometry}
\usepackage{amsmath}
\usepackage[thmmarks, amsmath, thref]{ntheorem}
\usepackage{amsfonts}
\usepackage{amssymb}
\usepackage{fancyhdr}
\pagestyle{fancy}
\usepackage{anysize}
\usepackage{graphicx}
\usepackage{float}
\usepackage{hyperref}
\usepackage{enumitem}
\newtheorem{theorem}{Theorem}
\theoremsymbol{\ensuremath{\blacksquare}}
\newtheorem*{solution}{Solución:}

\usepackage{xcolor}

% Margenes y formalidades

\textwidth 6.25in % Margen Derecho y ancho del texto
\textheight 8.5in
\oddsidemargin 0in % Margen Izquierdo
\headheight 0in % Largo del texto
\leftmargin 1in
\parindent 0em % Salto entre la indentación (Sangría)https://www.overleaf.com/4519127669sdpptkhvsbfn
\topmargin -0.5in % Margen Superior
\parskip 0.5ex % Salto entre parrafos
\baselineskip 1pt



\lhead{
\includegraphics[scale=0.09]{figuras/logoUSM.png}} % texto izquierda de la cabecera
%\chead{}                                      % texto centro de la cabecera
\rhead{Universidad Técnica Federico Santa María\\Ingeniería Comercial } % número de página a la derecha
\renewcommand{\headrulewidth}{0.4pt} % grosor de la línea de la cabecera
\renewcommand{\footrulewidth}{0.4pt} % grosor de la línea del pie

\begin{document}
\vspace{5cm}



\begin{center}
\begin{large}
\textbf{\\Gu\'ia Certamen 2}
\textbf{Econometría (ICS-294)}
\end{large}\\
\end{center}


\section{Preguntas Cortas}

\begin{enumerate}
    \item Se estima el siguiente modelo para explicar las ventas (\emph{SALES})
        en función del gasto en publicidad (\emph{ADV}):
        \[
          \widehat{\log(\text{SALES})}
          = 1.2 + 0.45\,\log(\text{ADV}) ,
        \]
        donde \texttt{SALES} se mide en millones de dólares y \texttt{ADV}
        en miles de dólares. Interprete el coeficiente \(\hat{\beta}_1 = 0.45\). ¿Qué concepto económico representa este coeficiente?\\

      {\color{blue}  \textbf{Solución.}
        \begin{itemize}
          \item En un modelo \emph{log-log}, el coeficiente \(0.45\)
                representa una \textbf{elasticidad}.
                Específicamente, indica que un aumento de \(1\%\) en el gasto
                en publicidad se asocia, en promedio, con un aumento de
                aproximadamente \(0.45\%\) en las ventas, manteniendo todo
                lo demás constante.
        \end{itemize}
        }
        
\item El salario promedio (\emph{salario}) en función de la experiencia laboral
        (\emph{exper}) está dado por:
        \[
          \widehat{\text{salario}}
          = 2.5 + 0.25\,exper - 0.005\,exper^2 .
        \]
        Determine el nivel de experiencia (\texttt{exper}) en el cual los
                retornos salariales a una unidad adicional de experiencia
                se anulan.
 \\
  {\color{blue} \textbf{Solución.}
        \begin{itemize}
          \item El efecto marginal de la experiencia es:
                \[
                  \frac{\partial\,\widehat{\text{salario}}}{\partial\,exper}
                  = 0.25 - 0.01\,exper.
                \]
          \item El retorno marginal se anula cuando:
                \[
                  0.25 - 0.01\,exper = 0
                  \quad\Rightarrow\quad
                  exper^\star = \frac{0.25}{0.01} = 25.
                \]
          \item Por tanto, a los \(25\) años de experiencia,
                el salario alcanza su máximo.
                A partir de ese punto, años adicionales de experiencia
                comienzan a disminuir el salario esperado.
        \end{itemize}}

\item ¿Por qué se asume la normalidad de las perturbaciones \(u\)
        en el modelo lineal clásico? Explique una razón   que justifiquen su importancia para realizar inferencia
        sobre los coeficientes estimados.
        \\
\\
{\color{blue} \textbf{Solución:} La normalidad garantiza que los estimadores de mínimos
                cuadrados ordinarios (MCO) tengan una distribución exacta
                normal en muestras finitas, permitiendo realizar inferencias
                precisas (tests de hipótesis, intervalos de confianza).}


\end{enumerate}

\section{Ejercicios}

\begin{enumerate}
    \item En una regresión lineal simple, el coeficiente asociado a la variable
        \(X_1\) fue estimado en:
        \[
          \hat\beta_1 = 2.7,
          \qquad s.e.(\hat\beta_1) = 0.9,
          \qquad n = 250.
        \]
        Considere un nivel de significancia del 5\%.

        \begin{enumerate}[label=\alph*)]
          \item Plantee las hipótesis nula y alternativa para evaluar si
                el coeficiente de la pendiete es significativamente distinto de cero.
      Calcule el estadístico \( t \) de prueba, compare su valor
                con el valor crítico correspondiente y determine si se
                rechaza la hipótesis nula.
          \item Construya un intervalo de confianza al 95\% para
                \( \beta_1 \) e interprete su significado.
                \item Suponga que desea contrastar si la variable dependiente presenta una relación creciente con respecto a la variable explicativa $X_1$.
Formule claramente las hipótesis nula y alternativa correspondientes, y evalúe si la evidencia empírica permite rechazar la hipótesis nula.
        \end{enumerate}

  {\color{blue}       \textbf{Solución.}
        \begin{enumerate}
          \item Hipótesis a testear:
                \[
                  H_0: \beta_1 = 0
                  \qquad\text{versus}\qquad
                  H_1: \beta_1 \neq 0.
                \]
       El estadístico \( t \) es:
                \[
                  t
                  = \frac{\hat\beta_1 - 0}{s.e.(\hat\beta_1)}
                  = \frac{2.7}{0.9}
                  = 3.0.
                \]
                

\begin{figure}[H]
    \centering
    \includegraphics[width=1.15\linewidth]{Figuras/t_248.png}
\end{figure}

           
                Dado que \(|3.0| > 1.97\), \textbf{rechazamos} \( H_0 \).
                Existe evidencia suficiente para afirmar que \( \beta_1 \) es
                significativamente distinto de cero al 5\% de significancia.
          \item El intervalo de confianza al 95\% es:
                \[  \hat\beta_1 \pm t_{2.5\%} \times s.e.(\hat\beta_1)
                  
                  \left[
                    2.7 - 1.97 \times 0.9;
                    2.7 + 1.97 \times 0.9
                  \right]
                  =
                  \left[
                    0.927; 4.473
                  \right].\]
            
                \textbf{Interpretación:} Con 95\% de confianza, se espera que
                el efecto parcial de \( X_1 \) sobre la variable dependiente
                esté entre 0.927 y 4.473 unidades. Es decir, un aumento en
                \( X_1 \) está asociado, en promedio, a un aumento positivo y
                significativo en la variable dependiente.


                \item \textbf{Hipótesis:}
\begin{itemize}
    \item Hipótesis nula (\( H_0 \)): \( \beta_1 =0  \)
    \item Hipótesis alternativa (\( H_1 \)): \( \beta_1 > 0 \)
\end{itemize}

Esto corresponde a una prueba unilateral a la derecha para el parámetro \( \beta_1 \).

Supongamos que al estimar el modelo se obtiene:
\[
\hat{\beta}_1 = 2{,}7 \quad \text{y} \quad \text{error estándar} = 0{,}8
\]

El estadístico de prueba es:
\[
 t = \frac{2{,}7}{0{,}9} = 3
\]

Con un nivel de significancia del 5\% y usando una prueba unilateral, se compara este valor con   $t_{5\%}=1.65$.

Como el estadístico t> 1.65, existe evidencia estadística suficiente para concluir que la variable dependiente es creciente respecto de \( x_1 \).
        \end{enumerate}}

\item Se tiene un modelo de regresión lineal de la forma:
$$y=\beta_0 + \beta_1 x_1+\beta_2x_2+\beta_3 x3 + u$$
los resultados de la regresión con 250 datos son los siguientes:
        \[
          \hat{\beta} = \begin{bmatrix} 3 \\ 9 \\ 5 \\2 \end{bmatrix},
        \]
y 
        \[
          (X'X)^{-1} =
          \begin{bmatrix}
1 & 0  & 0  &  0\\
0 & 1  & 0  & -1\\
0 & 0  & 2  & -1\\
0 & -1 & -1 &  1
\end{bmatrix},
          \quad \hat{\sigma}^2 = 2
        \]
\begin{enumerate}
    \item Plantee la hipótesis que testea que $\beta_1=4\beta_3$.

    \item ¿Se puede rechazar la hipótesis nula con la información que disponemos con una significancia del 5\%?

\end{enumerate}

{\color{blue}

\textbf{Solución}
\begin{enumerate}
    \item $$H_0 : \beta_1-4\beta_3 = 0$$
$$H_1 : \beta_1-4\beta_3 \neq 0$$
\item el estad\'istico t esta dado por:
$$t = \frac{\beta_1-4\beta_3}{s.e.(\beta_1-4\beta_3)}$$

\[
\text{Var}(\hat{\beta}) = 2 \cdot
\begin{bmatrix}
1 & 0  & 0  &  0\\
0 & 1  & 0  & -1\\
0 & 0  & 2  & -1\\
0 & -1 & -1 &  1
\end{bmatrix}
=
\begin{bmatrix}
2 & 0  & 0  &  0\\
0 & 2  & 0  & -2\\
0 & 0  & 4  & -2\\
0 & -2 & -2 &  2
\end{bmatrix}
\]



 $s.e.(\beta_1-4\beta_3)$ esta dado por:
$$s.e.(\beta_1-4\beta_3)=\sqrt{Var(\beta_1-4\beta_3)}=\sqrt{Var(\beta_1)+16Var(\beta_3)-8Cov(\beta_1,\beta_3)}$$
$$\sqrt{Var(\beta_1)+16Var(\beta_3)-8Cov(\beta_1,\beta_3)}=\sqrt{2+16\cdot 2-8 \cdot (-2)}=\sqrt{50}$$
por tanto 
$$t = \frac{9-4\cdot2}{7.07}=0,14$$
 No se rechaza $H_0$,  $\beta_1$ puede  ser $4\beta_3$. 


\end{enumerate}}

\item Considere la siguiente regresión lineal múltiple estimada con \(n = 75\) observaciones:

\[
\hat{y} \;=\; \hat{\beta}_{0} + \hat{\beta}_{1}x_{1} + \hat{\beta}_{2}x_{2} +\hat{\beta}_{3}x_{3} +\hat{\beta}_{4}x_{4} 
\]

El modelo arroja un coeficiente de determinación
\[
R^2 = 0.62.
\]

El investigador sospecha que las variables \(x_{3}\) y \(x_{4}\) no aportan capacidad explicativa.  
Para contrastar su hipótesis estima el siguiente modelo:

\[
\hat{y} \;=\;\hat{\gamma}_{0} + \hat{\gamma}_{1}x_{1} + \hat{\gamma}_{2}x_{2} ,
\qquad\text{obteniendo}\quad
R^2 = 0.55.
\]
Formule y realice un test de hipótesis que permita decidir, con un 5\% de significancia, si las variables \(x_{3}\) y \(x_{4}\) son relevantes para el modelo. ($F_{0.95,(2,70)}=3.13$)

{\color{blue}{\bf Solución:}

Queremos contrastar si las variables \( x_3 \) y \( x_4 \) son relevantes para explicar la variable dependiente. Para ello, se utiliza un test \( F \) para comparar el modelo completo con el modelo restringido.

\textbf{Hipótesis:}
\begin{itemize}
    \item \( H_0: \beta_3 = \beta_4 = 0 \) \quad (las variables \( x_3 \) y \( x_4 \) no son relevantes)
    \item \( H_1: \) Al menos uno de los coeficientes es distinto de cero.
\end{itemize}

\bigskip

\textbf{Estadístico \( F \):}

Dado:
\[
R^2_{\text{ur}} = 0.62, \quad R^2_{\text{r}} = 0.55
\]
\[
n = 75, \quad k_{\text{ur}} = 4, \quad k_{\text{r}} = 2
\]

El estadístico \( F \) se calcula como:
\[
F = \frac{(R^2_{\text{ur}} - R^2_{\text{r}})/j}{(1 - R^2_{\text{ur}})/(n - k_{\text{ur}}-1)}} = 
\frac{(0.62 - 0.55)/2}{(1 - 0.62)/(75 - 5)} = \frac{0.07/2}{0.38/70} = \frac{0.035}{0.00543} \approx 6.45
\]


Como el valor calculado \( F \approx 6.45 \) es mayor que el valor crítico \( F_{0.95,(2,70)} = 3.13 \), se \textbf{rechaza la hipótesis nula} al 5\% de significancia.
 Existe evidencia estadística suficiente para afirmar que las variables \( x_3 \) y \( x_4 \) son relevantes para el modelo.}

\end{enumerate}

\end{document}